\documentclass[
    language=korean,
    doctype=article,
    institution=none,
]{../../omnilatex}

\RequirePackage{config/document-settings}

\begin{document}

\maketitle

\section*{요약}
이 문서는 OmniLaTeX의 한국어 조판 기능을 보여줍니다. 한국어 소프트웨어 개발 문서를 예시로, 폰트 설정, 줄 간격 처리, 한영 혼합 조판 등의 기능을 소개합니다.

\tableofcontents

\section{서론}
OmniLaTeX은 LuaLaTeX와 luatexja 엔진을 기반으로 한국어 문서의 네이티브 조판을 제공합니다。

한국어 조판은 중국어, 일본어와 함께 CJK(중국어-일본어-한국어) 조판의 한 축을 형성합니다. 한국어의 조판 특성에는 한글의 음절 단위 결합, 한자 병용, 영문 및 숫자와의 혼합 사용 등이 포함됩니다. 이러한 특성을 올바르게 처리하는 것이 가독성 높은 한국어 문서 작성의 핵심입니다.

한국의 IT 산업에서 기술 문서의 품질은 제품의 경쟁력과 직결됩니다. 국내외 사용자를 동시에 고려한 문서 작성은 현대 소프트웨어 개발의 필수 요소가 되었습니다.

\section{한국어 조판 기능}

\subsection{기본 조판}
가을 하늘 공활한데 높고 구름 없이 밝은 달이 우리의 사랑을 비춰주네。
산에는 낙엽만 지는구나. 봄을 놓친 나비는 어디로 갈까。 한글은 세종대왕이 창제한 과학적인 문자로, 초성, 중성, 종성의 조합으로 음절을 구성합니다。

한국어 단락은 원칙적으로 첫 줄 들여쓰기로 시작합니다. 줄 간격은 일반적으로 글자 크기의 1.5배에서 2배로 설정됩니다. 문장 부호의 크기와 위치는 한국어 조판 규칙에 따라 자동으로 조정됩니다。

\subsection{한영 혼합 조판}
한국어와 영문이 혼합될 때 OmniLaTeX은 자동으로 적절한 간격을 삽입합니다。
예를 들어: LuaLaTeX 컴파일 엔진과 luatexja 패키지를 사용하여 고품질 한국어 조판을 구현할 수 있습니다。

한국어 기술 문서에서는 영문 용어의 사용이 불가피합니다. 예를 들어,「알고리즘의 시간 복잡도는 $O(n \log n)$ 이다」「RESTful API를 사용하여 데이터를 전송한다」와 같은 문장에서 한글과 영문 사이의 간격 처리가 중요합니다. OmniLaTeX은 이러한 간격을 자동으로 계산하여 일관된 조판 결과를 보장합니다。

\subsection{루비 주석}
OmniLaTeX은 루비 주석을 지원합니다：

\ruby{漢}{한}\ruby{字}{자}

루비 주석은 주로 한자의 음독을 표시하는 데 사용됩니다. 한국어 문서에서 한자는 법률 문서, 학술 논문, 전공 서적 등에서 여전히 널리 사용되며, 루비 주석을 통해 독자의 이해를 돕습니다。 또한 전문 용어의 축약형을 표기하는 용도로도 활용할 수 있습니다。

\subsection{수학식 조판}
한국어 학술 문서에서 수학식의 정확한 조판은 필수적입니다. 예를 들어, 피보나치 수열의 점화식은 $F(n) = F(n-1) + F(n-2)$으로 표현되며, 고윳값 분해의 수식은 다음과 같이 작성할 수 있습니다：

$$A = P \Lambda P^{-1}$$

OmniLaTeX은 \LaTeX{}의 수학 조판 기능을 한국어 문서에서도 완벽하게 지원합니다. 행 내 수식, 독립 수식, 행렬, 다행 정렬 등의 기능을 자유롭게 사용할 수 있습니다。

\section{한국어 기술 문서 작성 실무}

\subsection{API 문서}
소프트웨어 API 문서는 개발자가 가장 자주 참조하는 기술 문서 유형입니다. 한국어 API 문서에는 엔드포인트 URL, 요청 방법, 파라미터 정의, 응답 형식, 오류 코드 설명이 포함되어야 합니다。 OmniLaTeX의 표 환경과 코드 강조 기능을 활용하면 전문적인 수준의 API 문서를 쉽게 작성할 수 있습니다。

\subsection{논문 작성}
한국어 학술 논문은 초록, 서론, 본론, 결론의 네 부분 구성이 일반적입니다. 이공계 논문에서는 수식과 표의 정확한 조판이 필수적입니다。 한국어 논문은 학위 논문의 경우 교육부 학위논문 표지양식을 따르는 경우가 많으며, OmniLaTeX의 institution 옵션을 통해 이러한 형식 요구사항을 충족할 수 있습니다。

\subsection{오픈소스 프로젝트 문서}
한국의 오픈소스 생태계가 성장함에 따라, 한국어 오픈소스 프로젝트 문서의 수요도 증가하고 있습니다。 OmniLaTeX은 프로젝트 소개서, 기여 가이드, 아키텍처 설계 문서, 릴리스 노트 등 다양한 문서 유형의 조판에 사용할 수 있습니다. 통일된 조판 스타일은 프로젝트의 전문성과 커뮤니티 신뢰도를 높입니다。

\subsection{수식 예시}
한국어 수학 및 이공계 문서에서는 복잡한 수식의 조판이 자주 필요합니다。 예를 들어, 오일러의 등식은 $e^{i\pi} + 1 = 0$으로 표기되며, 가우스 적분은 다음과 같이 조판할 수 있습니다：
\[
\int_{-\infty}^{+\infty} e^{-x^2} \, dx = \sqrt{\pi}
\]

\section{결론}
OmniLaTeX은 한국어 학술 문서에 완전한 조판 솔루션을 제공합니다。 LuaLaTeX와 luatexja의 통합을 통해 루비 주석, 한영 혼합 간격 조정, 수학식 조판 등 한국어 특유의 조판 요구사항을 자동으로 처리합니다。 기술 문서부터 학술 논문까지, 폭넓은 장르의 한국어 문서를 고품질로 조판할 수 있습니다。

\end{document}
